% Generated draft: requires mathematical and source review.
\documentclass[11pt]{article}
\usepackage[margin=25mm]{geometry}
\usepackage{fontspec}
\setmainfont{Latin Modern Roman}
\usepackage{amsmath,amssymb,mathtools,unicode-math}
\setmathfont{Latin Modern Math}
\usepackage{enumitem,needspace,xurl,hyperref}
\setlength{\parindent}{0pt}
\setlength{\parskip}{5pt}
\setlength{\emergencystretch}{3em}
\newcommand{\N}{\mathbb{N}}
\newcommand{\Z}{\mathbb{Z}}
\newcommand{\Q}{\mathbb{Q}}
\newcommand{\R}{\mathbb{R}}
\newcommand{\C}{\mathbb{C}}
\newcommand{\F}{\mathbb{F}}
\newcommand{\A}{\mathbb{A}}
\newcommand{\PP}{\mathbb{P}}
\newcommand{\E}{\mathbb{E}}
\newcommand{\eps}{\varepsilon}
\newcommand{\dd}{\,\mathrm{d}}
\begin{document}
\section*{UnsolvedMath ID: 801030}
\textbf{Source problem code:} catalogue entry
\subsection{Definitions and mathematical statement}
A Turing degree is an equivalence class of subsets of \(\N\) under mutual Turing computability; write \(\mathbf a\leq_T\mathbf b\) for the induced ordering. The supplied record does not define degrees of categoricity or degrees of decidable categoricity. The following makes the question precise under the conventional presentation-based interpretation; its agreement with the source terminology requires review.

Work with countable structures in computable first-order languages, presented on computable domains contained in \(\N\). A presentation is computable if its atomic diagram is computable, and decidable if its elementary diagram (all first-order formulas with parameters that hold in the presentation) is computable. An isomorphism is computable in a degree \(\mathbf d\) if its graph is computable relative to a representative of \(\mathbf d\).

For a computable structure \(\mathcal A\), let \(C_{\mathcal A}(\mathbf d)\) mean that every computable presentation \(\mathcal B\cong\mathcal A\) admits an isomorphism \(\mathcal A\to\mathcal B\) computable in \(\mathbf d\). A degree \(\mathbf d\) is a degree of categoricity if there exists such an \(\mathcal A\) for which
\[
C_{\mathcal A}(\mathbf d)\quad\text{and}\quad
\forall\mathbf e\,\bigl(C_{\mathcal A}(\mathbf e)\Longrightarrow\mathbf d\leq_T\mathbf e\bigr).
\]
For a decidable structure \(\mathcal A\), let \(D_{\mathcal A}(\mathbf d)\) mean that every decidable presentation \(\mathcal B\cong\mathcal A\) admits an isomorphism \(\mathcal A\to\mathcal B\) computable in \(\mathbf d\). Under the interpretation used here, \(\mathbf d\) is a degree of decidable categoricity if there exists such a decidable \(\mathcal A\) for which
\[
D_{\mathcal A}(\mathbf d)\quad\text{and}\quad
\forall\mathbf e\,\bigl(D_{\mathcal A}(\mathbf e)\Longrightarrow\mathbf d\leq_T\mathbf e\bigr).
\]
In both definitions, least degree, rather than merely a minimal degree, is required, and the isomorphism may depend on the presentation.

The source states that every degree of decidable categoricity is a degree of categoricity. Its main question is whether
\[
\forall\mathbf d\;\bigl(\mathbf d\text{ is a degree of categoricity}
\Longrightarrow
\mathbf d\text{ is a degree of decidable categoricity}\bigr).
\]
The witnessing structures in the two definitions need not be the same.

The source further conjectures that this implication holds for degrees \(\mathbf d\geq_T\mathbf 0^{(\alpha)}\), where \(\alpha\) is a sufficiently large computable ordinal, and suggests \(\alpha=\omega\). Here \(\mathbf 0^{(\alpha)}\) denotes the degree of the transfinite iterated Turing jump of the empty set, using an effective ordinal notation; in particular,
\[
\mathbf 0^{(\omega)}=\deg_T\!\left(\bigoplus_{n\in\N}\emptyset^{(n)}\right).
\]
The explicit suggested case asks whether
\[
\forall\mathbf d\geq_T\mathbf 0^{(\omega)}\;\bigl(\mathbf d\text{ is a degree of categoricity}
\Longrightarrow
\mathbf d\text{ is a degree of decidable categoricity}\bigr).
\]
The phrase ``sufficiently large'' supplies no definite ordinal threshold or precise quantification over ordinal notations.

Finally, a d.c.e. degree is a Turing degree containing a set \(U\setminus V\), where \(U,V\subseteq\N\) are computably enumerable. The source asks what happens for d.c.e. degrees \(\mathbf d\). In the context of its main question, this asks whether every d.c.e. degree that is a degree of categoricity is also a degree of decidable categoricity; the wording does not specify whether a broader classification is intended.
\subsection{Short English statement}
Does every degree of categoricity also occur as a degree of decidable categoricity? The source conjectures an affirmative answer above a sufficiently high computable jump, suggesting \(\mathbf 0^{(\omega)}\), and separately asks about d.c.e. degrees. The definitions used here are an explicit conventional interpretation requiring source confirmation. The registry labels the problem open; this is supplied metadata, not an independently verified status assessment.
\subsection{Sources}
\begin{itemize}
\item Local UnsolvedMath registry, permanent ID \(801030\), source problem code \texttt{catalogue entry}, title ``Open Problems in Mathematical Logic --- Problem 42.'' The supplied record attributes the problem to Barmpalias, Gao et al., with proposed year 2026, and labels it open.
\item Supplied source list: Barmpalias, Gao et al., \emph{Open Problems in Mathematical Logic} (2026), Problem 42, \url{https://arxiv.org/abs/2608.26628}. The record reports access on 2026-10-16 and extraction by \texttt{arxiv\_latex}. These are supplied provenance claims; the source and access were not independently verified for this draft.
\item Supplied upstream ID: \(801030\); upstream revision: \texttt{d371ee917ac9647fbb06cd6fc529deec6f3dda82}. No separate authoritative reference defining the technical terminology is supplied.
\end{itemize}
\end{document}
